\documentclass[11pt]{article}

\usepackage[a4paper,margin=1in]{geometry}
\usepackage{amsmath,amssymb}
\usepackage[hidelinks]{hyperref}

\newcommand{\F}{\mathcal{F}}
\newcommand{\Mol}{\mathrm{Mol}}
\newcommand{\ext}{\mathrm{ext}}
\emergencystretch=3em

\title{Literature Status: The Extreme-Point Problem In arXiv:2302.13951 Is Answered By arXiv:2412.04312}
\author{Automatic math-discovery smoke test}
\date{June 29, 2026}

\begin{document}
\maketitle

\section*{Status}

This is a scoped \textbf{literature\_already\_answered} packet.  It records
that one open-problem signal in Aliaga--Pernecka--Smith,
\emph{Convex integrals of molecules in Lipschitz-free spaces},
arXiv:2302.13951 \cite{APS2023}, is answered by the later paper
Aliaga--Pernecka--Smith, \emph{A solution to the extreme point problem and
other applications of Choquet theory to Lipschitz-free spaces},
arXiv:2412.04312 \cite{APS2024}.

\section*{Source Open Problem Signal}

In Section 6 of arXiv:2302.13951, titled ``Applications to extremal
structure'', the source paper says that the main open problem in this direction
is the characterization of the extreme points of the unit ball of
\(\F(M)\).  More specifically, it identifies the unproved conjecture that all
extreme points must be elementary molecules.  The same paragraph notes that the
conjecture was known for proper metric spaces and subsets of \(\mathbb{R}\)-trees,
and then proves the claim for extreme points representable as convex integrals
of molecules.

\section*{Supporting Answer}

The later paper arXiv:2412.04312 explicitly frames the same problem as the
long-standing extreme point problem for Lipschitz-free spaces.  Its Section 4,
``The extreme point problem'', states that the problem is to determine whether
all extreme points must be molecules, and then says that Theorem
4.1 provides a definitive answer.

\begin{quote}
\textbf{Theorem 4.1 in arXiv:2412.04312.}
Let \(M\) be a complete metric space. Then
\(\ext B_{\F(M)}\subset \Mol(M)\). Specifically, the extreme points of
\(B_{\F(M)}\) are exactly the molecules \(m_{xy}\), \(x\ne y\), such that
\[
  d(x,p)+d(p,y)>d(x,y)
  \qquad(p\in M\setminus\{x,y\}).
\]
\end{quote}

\section*{Conclusion}

This answers the extreme-point conjecture/open problem mentioned in Section 6
of arXiv:2302.13951, for the same complete-metric-space setting used in the
source paper.

\section*{Scope Limitations}

This packet is deliberately narrow.  It does not claim that arXiv:2412.04312
answers the source paper's explicit Question in Section 3 about rewriting a
convex series of molecules so that no first coordinate is a second coordinate.
It also does not claim to answer the explicit Question in Section 4 asking
whether the isometric-copy hypothesis in Theorem 4.1 of arXiv:2302.13951 can
be replaced by a bi-Lipschitz-copy hypothesis.  In fact, arXiv:2412.04312
recasts a closely related curve-fragment question as an open question.

\section*{Evidence Checked}

The local source TeX for arXiv:2302.13951 was inspected around the source
open-problem signal and the two explicit questions.  The local source TeX for
arXiv:2412.04312 was inspected around its abstract, introduction, and Theorem
\texttt{th:extreme}.  A web check on June 29, 2026 found the arXiv pages for
both the source paper and the supporting paper.  The run indexes had an
adjacent literature-status packet for another source paper and arXiv:2412.04312,
but no exact packet for arXiv:2302.13951.

\begin{thebibliography}{9}
\bibitem{APS2023}
Ram\'on J. Aliaga, Eva Perneck\'a, and Richard J. Smith,
\emph{Convex integrals of molecules in Lipschitz-free spaces},
arXiv:2302.13951, 2023.

\bibitem{APS2024}
Ram\'on J. Aliaga, Eva Perneck\'a, and Richard J. Smith,
\emph{A solution to the extreme point problem and other applications of
Choquet theory to Lipschitz-free spaces},
arXiv:2412.04312, 2024.
\end{thebibliography}

\end{document}
